% !TEX program = pdflatex
% Lettre de motivation — Siemens Industry Software SARL · Software Engineer · Offre 504804
\documentclass[11pt,a4paper]{letter}

\usepackage[french]{babel}
\usepackage[utf8]{inputenc}
\usepackage[T1]{fontenc}
\usepackage{lmodern}
\usepackage[margin=2.2cm]{geometry}
\usepackage{hyperref}
\usepackage{microtype}
\usepackage{xcolor}

\hypersetup{
  colorlinks=true,
  urlcolor=blue,
  pdfauthor={Hamouda Chekir},
  pdftitle={Lettre de motivation — Hamouda Chekir — Siemens EDA Software Engineer}
}

\address{%
  Hamouda Chekir\\
  Tunis, Tunisie\\
  +216 55 913 832\\
  \href{mailto:hamoudachkir@yahoo.fr}{hamoudachkir@yahoo.fr}%
}

\date{\today}
\signature{Hamouda Chekir}

\begin{document}

\begin{letter}{%
  Siemens Industry Software SARL\\
  Service Recrutement\\
  Tunis, Tunisie\\
  Offre n\textsuperscript{o} 504804 — Software Engineer%
}

\opening{Madame, Monsieur,}

Actuellement en fin de cycle d'ingénieur à l'ESPRIT (spécialité TWIN), je candidate au poste de \textbf{Software Engineer} au sein de la division Hardware-Assisted Verification (offre 504804), pour rejoindre l'équipe R\&D de prototypage FPGA autour de la plateforme Veloce proFPGA.

Votre offre recherche un profil débutant capable de concevoir, modifier et déboguer du logiciel de qualité, en C++ et avec une culture d'algorithmes et de POO, tout en évoluant sous encadrement senior. C'est exactement le cadre dans lequel je souhaite démarrer : apprendre vite, livrer proprement, et renforcer la qualité fonctionnelle des produits sur les plateformes cibles.

Chez \textbf{Talan}, dans le cadre de mon projet de fin d'études, j'ai contribué à NextHire, une plateforme de \textbf{16 microservices}. J'y ai travaillé l'intégration de composants, le debugging de parcours temps réel (WebRTC, contrôles d'intégrité) et la stabilisation des releases — des situations où traquer une régression, comprendre un mode de panne et livrer une correction fiable importaient autant qu'écrire une nouvelle fonctionnalité. Au quotidien, je m'appuie sur \textbf{Linux, Python, Bash et Git}, des réflexes directement utiles dans un environnement EDA (scripting, automatisation, collaboration multi-équipes).

Ma base en \textbf{C/C++, structures de données et POO} vient de ma formation d'ingénieur et de ma pratique personnelle. Je n'ai pas encore d'expérience production en HDL (VHDL, Verilog, SystemVerilog) ni sur Vivado/Quartus/STA : je le dis clairement. En revanche, je suis motivé à monter rapidement sur ces briques, sous mentorat, pour contribuer au logiciel de validation bas niveau et à l'intégration système autour du prototypage FPGA. Tunis, CDI, travail sur site : je suis disponible et engagé sur ce format.

Anglais courant (écrit et oral), habitué aux revues de code et aux équipes transverses, je souhaite mettre mon sérieux technique et ma curiosité au service de Siemens EDA.

Je reste à votre disposition pour un entretien.

\vspace{0.5em}
\closing{Je vous prie d'agréer, Madame, Monsieur, l'expression de mes salutations distinguées.}

\end{letter}
\end{document}
